% Conditional reduction and diagnostic candidates. Neither candidate is proved.
\section{A clipped forward comparison for a weaker induction input}
Let $A=F_X$, $B=F_Y$, with weights $a,b$, and put
\[
 D_n=\max(A_n-A_{n-1},0)\quad(n\ge2),\qquad H=E_B A.
\]
Here $D_0=D_1=0$. The following is an \emph{unproved}
generation-dependent comparison:
\[
 \frac{\min(H_n,D_n)}{D_n}\le
 \frac{\min(H_{n+1},D_{n+1})}{D_{n+1}}
 \quad\left(2\le n<\min(a,b),\ D_nD_{n+1}>0\right).
                                                        \tag{R}
\]
All ratios can be cross multiplied. Full, unclipped ratio
monotonicity is false; subtracting a single unit does not
repair it either. The clipping and the stated range must
not be silently removed.

\begin{proposition}[Conditional coefficient-tilt consequence]
If (R) holds for a given pair of increasing nonrectangular
arms, then the coefficients of $A(z/2)B(z/2)$ are unimodal.
Equivalently $2^j(q_Xq_Y)_j$ is unimodal.
This is coefficient tilt, not the weighted-positive-difference
property needed to finish the original conjecture.
\end{proposition}
\begin{proof}
The exact identity
\[
 (2z-1)AB=E_B A-(1-z)A                              \tag{1}
\]
shows that, while $A$ is rising, its rank-$n$ coefficient
is $H_n-D_n$. Once $H_n\ge D_n>0$, (R) forces
$H_{n+1}\ge D_{n+1}$ whenever the latter denominator is
positive. If it vanishes, initial interval support of the
ordinary positive differences shows that $A$ has finished
rising; the coefficient in (1) is then nonnegative forever.

Only ranks through $\min(a,b)$ require this argument.
Orient $a\le b$. For $n>a$, every nonzero summand in
$[z^n]AB=\sum_{i=0}^{a}A_iB_{n-i}$ uses $n-i\ge1$.
The elementary actual-ideal growth bound
$B_r\le2B_{r-1}$ therefore gives
$[z^n]AB\le2[z^{n-1}]AB$ directly. The initial ranks
are $1,2,5$ in $AB$, so its tilt starts with a plateau
and a rise. Identity (1), (R) and the terminal growth bound
exclude every later decrease followed by a strict rise.
\end{proof}

\section{A second, still open, coupled comparison}
For a legal append $X=Zt$ orient $a=\mathrm{wt}X\ge b$,
put $p=q_Zq_Y$, and let $m$ be its last mode and $s$ the
first mode of the new ordinary product. Let
$u=(\Delta p)^+$ and let $\ell$ be its first weighted
mode, that is, a first mode of $2^j u_j$.
In the crossing range $m<n<s$, put $k=n-t$ and
\[
 w_n=p_n-p_k,\qquad v_n=p_{n-1}-p_n.
\]
The following sufficient resource is also \emph{unproved}:
\[
 u_k+v_n\le2(u_{k+1}+v_{n+1})                       \tag{S}
\]
when $0\le k<\ell$, $w_n>0$,
$2w_n<w_{n-1}$, and $2p_n\ge p_{n-1}$.
Broader versions of (S), with the last condition removed,
are false, including examples with $k=0$.

The reason (S) is useful is exact and does not rely on any
finite diagnostic. In this range
\[
 w_{n-1}-w_n=u_k+v_n,\qquad
 w_n-w_{n+1}=u_{k+1}+v_{n+1}.
\]
Thus a strict weighted fall propagates to the next edge
under (S). If instead the tail $2^n p_n$ has already
strictly fallen, a proved coefficient-tilt assertion would
control it while the outgoing head is still increasing.
For $k\ge\ell$, the geometric tail bound of the old
weighted positive differences handles the window directly.
These observations motivate (R) and (S), but neither is
asserted to hold for all actual arms.

\texttt{ClippedForwardRatio.lean} checks (1), the unbounded
clipped-threshold implication, weighted finite-sum growth,
the short-factor index arithmetic, and strict propagation
under the two-loss resource. This does not verify (R),
(S), or a complete ordinary-product induction.

\paragraph{A valid transport fact, and the remaining source.}
Suppose $D$ is positive on an interval and $H$ vanishes to
its left. If $\min(H/D,1)$ is nondecreasing there and
$H\ge0$ everywhere, convolution of both $H,D$ by the same
interval preserves the clipped ratio order. For each fixed
$0<\lambda<1$, the input $H-\lambda D$ changes sign at
most once, from nonpositive to nonnegative. A uniform window
preserves this sign pattern: while it straddles the change,
it loses a nonpositive term and gains a nonnegative one.
Outside that range all its terms have one sign. A decrease
of the output clipped ratio would contradict this fact for
a $\lambda$ strictly between its two values.

This transport fact alone does not prove (R). On deleting
the last part of the $A$ arm, before the ordinary terminal
splice the actual pair changes by
\[
 H_{\rm new}=E_B+z[t]H_{\rm old},\qquad
 D_{\rm new}=z^2[t-1]+z[t]D_{\rm old}.
\]
The different source terms must be controlled together;
discarding them or asserting that addition preserves clipped
ratio order would be an unjustified closure step. These
identities isolate the still missing generation argument.

\section{Reduction of (R) to a single-arm shifted cumulative comparison}
For a single increasing nonrectangular arm $A$, write
$G_j=\sum_{i\le j}A_i$, zero when $j<0$. The following
single-arm assertion is still \emph{unproved}:
\[
 G_{n+1-r}D_n\ge G_{n-r}D_{n+1}
 \quad\left(r\ge3,\ D_nD_{n+1}>0,\ n\ge2\right).
                                                        \tag{C}
\]
Unlike (R), this assertion has no clipping. Broad and targeted
finite tests do not constitute its proof.

\begin{proposition}[Conditional single-arm reduction]
If (C) holds for every increasing arm, then (R) holds for
every pair of such nonrectangular arms. Consequently their
ordinary product has unimodal factor-two coefficient tilt.
\end{proposition}
\begin{proof}
At a pair of ranks where $B$ is still strictly rising,
the established single-arm deficit mode ordering gives
$\Delta E_B(j)\ge0$ throughout the prefix in question.
Since $E_B(j)=0$ for $j\le2$, summation of its differences
gives the exact coefficient identity
\[
 H_n=(E_B A)_n
       =\sum_{j=3}^{n}\Delta E_B(j)\,G_{n-j}.
                                                        \tag{2}
\]
Apply (C) to every $j$ in this sum. The new term at $j=n+1$
has zero old coefficient and a nonnegative new coefficient.
Multiplying by the nonnegative $\Delta E_B(j)$ and summing
therefore proves the \emph{unclipped} ratio comparison
$H_{n+1}D_n\ge H_nD_{n+1}$ in this range. Clipping preserves
this comparison at threshold one.

If $B$ has already ceased strictly rising at rank $n+1$,
its positive-difference initial support implies it never
strictly rises again. Use the symmetric version of (1):
\[
 (2z-1)AB=E_A B-(1-z)B.
\]
Its coefficients at this and all later ranks are nonnegative,
because $E_A B\ge0$ and $B_j-B_{j-1}\le0$ there.
Thus $H_{n+1}\ge D_{n+1}$ whenever $D_{n+1}>0$, and
the new clipped ratio is one. The same argument handles
leaving the rising support of $A$ as in the earlier proof.
These cases establish (R).
\end{proof}

\texttt{weighted\_ratio\_sum} in
\texttt{ClippedForwardRatio.lean} checks the unbounded
weighted-sum comparison used in (2). The ordinary deficit
mode input and coefficient summation are the written proof.
No assertion of (C), (R), (S), or the general original
conjecture is made here.
